\chapter{可靠性门控的类条件时序自适应方法设计}

\section{方法设计动机}
第四章的实验结果表明，现有领域自适应基线虽然已经能够在目标域上取得明显收益，但仍存在三个尚未解决的问题。首先，单纯的全局分布对齐难以显式保证跨域同类样本聚合、异类样本分离，因此当不同故障类别在目标域中边界接近时，模型仍可能出现较严重的类别混叠。其次，条件对齐方法虽然能够利用类别预测信息，但目标域伪标签本身存在噪声，若直接以其驱动对齐，容易把错误类别关系进一步强化。最后，从当前结果看，不同方法在场景之间的波动仍然较大，这意味着改进方法除了追求更高上限，还需要兼顾训练稳定性和目标域监督的可靠性。

基于上述问题，本文提出一种可靠性门控的类条件时序自适应方法（Reliability-Gated Class-Conditional Temporal Adaptation, RCTA）。该方法的核心思想是：利用教师网络生成更平滑的目标域预测，以多因素可靠性评分筛选可用伪标签，并结合类别原型约束与一致性正则，把全局域对齐和类条件迁移有机结合起来，从而提高跨工况故障诊断中的稳定迁移能力。

\section{方法总体框架}
本文提出的改进方法记为 RCTA。整体框架仍由特征提取器、分类器和域对齐模块三部分组成，其中 FCN 作为共享时序特征提取骨干，以保持与第四章基线比较的一致性。与现有基线相比，RCTA 的主要新增内容包括以下四项：

\begin{enumerate}
\item 教师网络伪标签估计：利用指数滑动平均教师网络生成更稳定的目标域预测；
\item 可靠性门控机制：综合置信度、熵、一致性和原型相似度筛选可信目标样本；
\item 类别原型建模：维护源域与目标域的类别原型，为类条件迁移提供参照；
\item 类条件与一致性约束：同时强化原型吸引、类别分离和时序一致性。
\end{enumerate}

在训练过程中，RCTA 保留现有全局域对齐损失作为基础约束，例如 CDAN 的条件对抗损失或 DANN 的域对抗损失；与此同时，通过可靠目标监督、双域原型和一致性正则进一步增强跨域同类聚合，从而缓解“全局对齐有效但类别边界不清晰”的问题。

\section{类别原型建模}
设特征提取器为 $\phi(\cdot)$，分类器为 $h(\cdot)$。对于一个源域小批量 $\mathcal{B}_s=\{(x_i^s,y_i^s)\}_{i=1}^{m_s}$，记其特征表示为
\begin{equation}
f_i^s=\phi(x_i^s).
\end{equation}
对于第 $k$ 类故障，本文以该类别样本在特征空间中的均值作为小批量原型：
\begin{equation}
\tilde{c}_k=\frac{1}{|\mathcal{B}_s^k|}\sum_{i:y_i^s=k} f_i^s,
\end{equation}
其中 $\mathcal{B}_s^k$ 表示当前批中属于第 $k$ 类的样本集合。为了减小单个小批量波动对类别中心的影响，本文采用指数滑动平均方式更新全局原型：
\begin{equation}
c_k \leftarrow \mu c_k + (1-\mu)\tilde{c}_k,
\end{equation}
其中 $\mu \in [0,1)$ 为动量系数。这样得到的源域类别原型一方面能够反映各故障类别的稳定中心，另一方面又能随训练过程逐步更新，作为目标域特征对齐的参照。

\section{目标域可靠伪标签筛选}
对于目标域小批量 $\mathcal{B}_t=\{x_j^t\}_{j=1}^{m_t}$，记其预测概率为
\begin{equation}
p_j^{t,\mathrm{tea}}=\mathrm{softmax}\!\left(\frac{h_{\mathrm{tea}}(\phi_{\mathrm{tea}}(x_j^t))}{T_{\mathrm{tea}}}\right),
\end{equation}
伪标签定义为最大概率对应的类别：
\begin{equation}
\hat{y}_j^t=\arg\max p_j^{t,\mathrm{tea}}.
\end{equation}
同时记预测置信度为
\begin{equation}
q_j^t=\max p_j^{t,\mathrm{tea}}.
\end{equation}
为了降低错误伪标签带来的负迁移风险，本文不直接让所有目标域样本参与类条件训练，而是为每个样本构造综合可靠性评分。除最大预测概率外，还考虑预测熵、学生网络与教师网络之间的一致性，以及目标样本与参考原型之间的相似度。记综合可靠性分数为
\begin{equation}
r_j^t=\alpha_1 q_j^t+\alpha_2(1-e_j^t)+\alpha_3 u_j^t+\alpha_4 s_j^t,
\end{equation}
其中 $e_j^t$ 表示归一化预测熵，$u_j^t$ 表示学生网络对伪标签的响应强度，$s_j^t$ 表示样本特征与对应类别参考原型之间的相似度，$\alpha_1,\alpha_2,\alpha_3,\alpha_4$ 为非负权重且满足和为 1。进一步通过阈值 $\tau$ 选出高可靠目标子集
\begin{equation}
\Omega_t=\{j \mid r_j^t > \tau\}.
\end{equation}
随着训练进行，阈值和对应损失权重采用渐进式策略，使模型先依赖源域监督和全局对齐获得相对稳定的表示，再逐步利用高可靠目标样本开展更细粒度的类条件迁移。

\section{类条件对齐与一致性约束}
\subsection{原型吸引损失}
对于被选中的高置信目标域样本，本文希望其特征尽量接近对应类别的源域原型，因此定义原型一致性损失为
\begin{equation}
\mathcal{L}_{proto}=
\frac{1}{|\Omega_t|}
\sum_{j\in\Omega_t}
\left\|
\phi(x_j^t)-c_{\hat{y}_j^t}
\right\|_2^2.
\end{equation}
该项约束的作用是直接缩小目标域特征与源域同类中心的距离，从而增强跨域同类样本的聚合性。

\subsection{原型分离约束}
除增强同类样本聚合外，本文还希望不同类别的参考原型保持足够分离。记参考原型集合为 $\{c_k\}_{k=1}^{K}$，则原型分离损失定义为
\begin{equation}
\mathcal{L}_{sep}=
\frac{1}{K(K-1)}\sum_{k\neq l}
\max\left(0,\cos(c_k,c_l)-m\right),
\end{equation}
其中 $m$ 为原型间余弦相似度的分离边界。该项约束能够抑制不同类别中心在特征空间中过度靠近，从而减少类间混叠。

\subsection{一致性监督项}
在高可靠目标样本基础上，本文进一步约束学生网络预测与教师网络预测保持一致。设学生网络预测概率为 $p_j^{t,\mathrm{stu}}$，则一致性损失可写为
\begin{equation}
\mathcal{L}_{cons}=
\frac{1}{|\Omega_t|}\sum_{j\in\Omega_t}
\mathrm{KL}\!\left(p_j^{t,\mathrm{tea}} \,\|\, p_j^{t,\mathrm{stu}}\right).
\end{equation}
该项的作用是进一步稳定目标域判别边界，并利用教师网络提供的平滑监督减轻伪标签噪声对训练的扰动。

\section{整体优化目标与训练策略}
在保留现有全局域对齐损失 $\mathcal{L}_{global}$ 的基础上，RCTA 的总损失函数可写为
\begin{equation}
\mathcal{L}=
\mathcal{L}_{cls}^s
+\lambda_g\mathcal{L}_{global}
+\lambda_p\mathcal{L}_{proto}
+\lambda_s\mathcal{L}_{sep}
+\lambda_c\mathcal{L}_{cons},
\end{equation}
其中 $\mathcal{L}_{cls}^s$ 表示源域监督分类损失，$\lambda_g$、$\lambda_p$、$\lambda_s$ 和 $\lambda_c$ 分别为各损失项的权重。训练策略上，本文按照以下方式组织优化：

\begin{enumerate}
\item 训练初期主要依赖源域分类损失和全局域对齐损失，先获得较稳定的跨域共享表示；
\item 随着训练推进，利用教师网络输出和可靠性门控机制筛选可信目标样本；
\item 在每轮迭代中更新类别原型，并通过原型吸引、原型分离和一致性监督共同约束目标域表示；
\item 最终在统一实验场景中评估该方法相对基线的目标域性能收益与稳定性变化。
\end{enumerate}

该设计与现有基线方法共享统一骨干网络和基本训练框架，便于在相同实验条件下分析新增机制的有效性。

\section{实验验证方案}
\subsection{与基线方法的总体比较}
实验验证首先在第四章已经完成的六组代表性单源场景上考察 RCTA 相对于 Source Only、DAN、DANN 和 CDAN 的增益情况；在此基础上，再扩展到全量场景的系统比较中进行补充说明。

\subsection{消融实验设计}
为了明确各模块贡献，本文设计以下消融版本：
\begin{enumerate}
\item 仅保留全局域对齐损失的基线版本；
\item 全局域对齐 + 原型吸引损失；
\item 全局域对齐 + 原型吸引 + 原型分离约束；
\item 全局域对齐 + 原型吸引 + 一致性监督；
\item 完整版本，即同时加入可靠性门控、原型约束和一致性正则。
\end{enumerate}
通过上述逐步叠加方式，可以分析每个模块对目标域性能、训练稳定性和类别区分能力的具体影响。

\subsection{重点分析指标}
在验证过程中，除了目标域准确率和平衡准确率外，还将重点关注以下现象：其一，源域准确率是否因为额外对齐约束而出现明显下降；其二，混淆矩阵中易混类别是否得到改善；其三，t-SNE 可视化中跨域同类样本是否比现有基线更加聚合；其四，所提方法在不同迁移方向上的收益是否更加稳定。

\section{本章小结}
本章在第四章实验结果分析的基础上，进一步给出了一个更具体的创新方向，即可靠性门控的类条件时序自适应方法。该方法尝试把教师网络伪标签、可靠性筛选、类别原型和一致性正则结合起来，以提升跨工况故障诊断中的类条件迁移能力。相较于仅停留在原则层面的改进设想，这一方案已经明确了核心机制、损失函数和验证路径，也为后续实验补充提供了清晰主线。
